\section{Use XOR as polynomial subtraction}
The checksum can be understood as a remainder calculation whose coefficients are
bits. A coefficient is zero or one. Adding two coefficient vectors uses XOR: a term
present twice cancels. There are no carries between adjacent coefficient positions,
and subtraction has the same effect as addition. For example, binary 101 XOR 011 is
110. Treating that operation as ordinary integer subtraction would instead produce
010, which is a different calculation.

A shift moves the coefficient positions. In the reflected implementation, inspect
the low bit before shifting right. If it is zero, the shifted state is the next state.
If it is one, XOR the shifted state with \texttt{edb88320}. That feedback cancels
the outgoing coefficient while accounting for the chosen divisor. For a 32-bit state
\(r\), one reflected step is
\[
 F(r)=\begin{cases}
 r\mathbin{\gg}1,&r\mathbin{\&}1=0,\\
 (r\mathbin{\gg}1)\mathbin{\oplus}\texttt{0xedb88320},&r\mathbin{\&}1=1.
 \end{cases}
\]
Here the right shift inserts zeroes. The notation describes a bit-vector operation;
it is not signed division with rounding. The mask and shifted state both fit 32 bits,
so the result fits 32 bits as well.

To consume an octet \(b\), XOR its unsigned value into the low byte of the state,
then apply \(F\) eight times. No byte is skipped because it is zero, and no zero
byte terminates the message. A byte count supplied by the caller determines the extent.
This matters for files, where zeroes are ordinary data rather than string terminators.
